\documentclass[conference,10pt]{IEEEtran}
\usepackage{cite}
\usepackage{amsmath,amssymb,amsfonts}
\usepackage{algorithmic}
\usepackage{graphicx}
\usepackage{textcomp}
\usepackage{xcolor}
\usepackage{booktabs}
\usepackage{microtype}
\usepackage{url}
\usepackage{hyperref}

\hypersetup{
    colorlinks=true,
    linkcolor=blue,
    filecolor=magenta,      
    urlcolor=cyan,
    citecolor=blue,
}

\begin{document}

\title{AdaptiveVec: A Density- and Dimension-Aware Proximity Graph Index for Resource-Constrained Vector Retrieval}

\author{
\IEEEauthorblockN{Aradhya Shinde}
\IEEEauthorblockA{\textit{Department of Computer Engineering} \\
\textit{Systems \& Machine Learning Research Laboratory}\\
Pune, India \\
aradhyashinde2330@gmail.com}
}

\maketitle

\begin{abstract}
Hierarchical Navigable Small World (HNSW) proximity graphs underpin modern approximate nearest neighbor search (ANNS) in production vector databases. However, canonical HNSW enforces rigid, globally uniform hyper-parameters ($M=16, efConstruction=200$) across heterogeneous high-dimensional manifolds. In dense, low-intrinsic-dimensionality subspaces, this uniform allocation synthesizes redundant proximity links, squandering memory and memory bus bandwidth with zero recall utility. Conversely, sparse high-dimensional boundary regions suffer from capacity starvation and topological disconnects, exacerbated by the hubness phenomenon. We introduce \textbf{AdaptiveVec}, a lightweight, manifold-adaptive proximity graph architecture tailored for resource-constrained commodity hardware ($100\text{K}$--$1\text{M}$ vectors on single-node laptops and edge instances). AdaptiveVec extracts online geometric signals---\textbf{Local Intrinsic Dimensionality (LID)} via maximum likelihood estimation and \textbf{Local Density ($D$)}---directly from standard greedy routing descent paths at under $0.8\%$ computational overhead, completely eliminating offline pre-clustering passes. It introduces: (1) a \textit{Layer-Decoupled Dynamic Allocation Policy} that scales per-node degree while compressing higher-layer express links; (2) \textit{Streaming Online Welford Tracking} with exponential decay; (3) a \textit{Hubness-Aware In-Degree Regulation Heuristic} that penalizes high-degree bottleneck vertices; (4) \textit{Distance Stagnation Early Exit} to truncate futile search hops; and (5) \textit{Asymmetric INT8 Scalar Quantization (SQ8)} with float32 distance re-ranking. AdaptiveVec delivers a \textbf{7.4\% reduction in graph edges}, up to \textbf{50.3\% higher query throughput (7,075.3 QPS vs. 4,708.1 QPS baseline)} via the distance stagnation early-exit mechanism, and \textbf{62.4\% index memory reduction (22.5MB vs. 59.9MB)} via asymmetric INT8 scalar quantization, at a measured Recall@10 of \textbf{0.9745} (throughput-optimized configuration) and \textbf{0.9594} (memory-optimized configuration) against a 0.9913 baseline, evaluated on a canonical SIFT-100K subset (Texmex IRISA) on commodity hardware (Intel Core 5 210H, 8 cores/12 threads, AVX2/FMA, 16GB RAM).
\end{abstract}

\begin{IEEEkeywords}
Approximate Nearest Neighbor Search (ANNS), Hierarchical Navigable Small World (HNSW), Local Intrinsic Dimensionality (LID), Hubness Phenomenon, Scalar Quantization, Edge Computing.
\end{IEEEkeywords}

\section{Introduction \& Motivation}
Dense vector embeddings generated by deep representation models underpin modern information retrieval, semantic recommendation systems, and Retrieval-Augmented Generation (RAG) pipelines. In these operational frameworks, high-throughput Approximate Nearest Neighbor Search (ANNS) over million-scale vector datasets is a non-negotiable prerequisite. Among available indexing paradigms, \textbf{Hierarchical Navigable Small World (HNSW)} graphs \cite{malkov2020hnsw} represent the gold standard in production systems (FAISS, Milvus, Qdrant, pgvector), achieving empirical logarithmic search complexity $\mathcal{O}(\log N)$ and recall exceeding 98\%.

\subsection{The Pathology of Uniform Allocation}
Despite widespread adoption, canonical HNSW suffers from a fundamental architectural limitation: \textbf{static, globally uniform parameterization}. Canonical implementations configure a single edge budget quota $M$ (e.g., $M=16$) and a single construction beam width $efConstruction$ (e.g., $efC=200$) applied identically across the entire dataset.

Real-world vector representations, however, are profoundly non-homogeneous. Vectors do not populate Euclidean space $\mathbb{R}^D$ uniformly; rather, they concentrate along complex, non-linear sub-manifolds characterized by localized variations in \textbf{Local Intrinsic Dimensionality (LID)} and \textbf{Local Density ($D$)}:
\begin{enumerate}
    \item \textbf{Redundant Edge Bloat:} In dense, low-LID subspaces (e.g., concentrated semantic clusters), vectors reside on low-dimensional hyperplanes. Maintaining 16 or 32 bidirectional edges generates severe structural redundancy: multiple parallel edges connect essentially co-linear neighbors, wasting DRAM and memory bandwidth while conferring zero routing advantage.
    \item \textbf{Hubness \& Congestion:} In sparse high-dimensional regions, uniform edge quotas starvations occur. Under the \textbf{Hubness Phenomenon} \cite{radovanovic2010hubs}, central nodes accumulate an exorbitant number of routing paths, generating edge-thrashing and degrading tail latency (P95/P99 SLAs).
    \item \textbf{The Commodity Hardware Barrier:} Enterprise vector search literature often evaluates on multi-socket servers with 128GB--512GB of RAM. In contrast, thousands of edge deployments and developer machines operate with strictly constrained memory and compute budgets ($100\text{K}$--$1\text{M}$ vectors on single-node laptops or cost-capped cloud virtual machines).
\end{enumerate}

To address these challenges, we present \textbf{AdaptiveVec}, a proximity graph architecture that dynamically adapts its topological budget based on online manifold geometry.

\section{Mathematical Formulation \& Online Signals}
Let $\mathcal{X} = \{x_1, x_2, \dots, x_N\} \subset \mathbb{R}^D$ denote a corpus of $N$ vectors embedded in an ambient space of dimension $D$, equipped with a metric distance function $d(x, y)$ (e.g., Euclidean $L_2$ distance).

\subsection{Local Intrinsic Dimensionality (LID) Estimation}
Given a local neighborhood of $k$ nearest candidate distances sorted in ascending order $d(x, v_1) \le d(x, v_2) \le \dots \le d(x, v_k)$, the Maximum Likelihood Estimator (MLE) of LID \cite{levina2005mle, amsaleg2015lid} is formulated as:
\begin{equation}
\widehat{\text{LID}}(x) = -\left[ \frac{1}{k - 1} \sum_{i=1}^{k-1} \ln \left( \frac{d(x, v_i)}{d(x, v_k)} \right) \right]^{-1}
\end{equation}
Where $d(x, v_k)$ serves as the maximum radius of the local exploratory ball. A small $\widehat{\text{LID}}(x)$ indicates a low-dimensional manifold (e.g., $2.0 \le \widehat{\text{LID}} \le 8.0$), whereas points in sparse high-dimensional clouds yield elevated values ($\widehat{\text{LID}} \ge 24.0$).

\subsection{Local Density Formulation}
To capture localized spacing independent of dimensional expansion rate, Local Density $D(x)$ is computed as the mean distance to the $k$ nearest discovered neighbors:
\begin{equation}
D(x) = \frac{1}{k} \sum_{i=1}^k d(x, v_i)
\end{equation}

\subsection{Online Streaming Welford Tracking}
To avoid multi-pass corpus pre-computation, AdaptiveVec tracks running statistical moments online using Welford's algorithm \cite{welford1962note}. For each inserted vector $x_n$, running mean $\bar{M}_n$ and squared variance accumulator $S_n$ are updated in $\mathcal{O}(1)$ time:
\begin{align}
\delta &= z_n - \bar{M}_{n-1} \\
\bar{M}_n &= \bar{M}_{n-1} + \frac{\delta}{n} \\
S_n &= S_{n-1} + \delta \cdot (z_n - \bar{M}_n)
\end{align}
The normalized difficulty score $S(x)$ combines standardized z-scores:
\begin{equation}
S(x) = \text{clip}\left( \alpha \cdot z_{\text{LID}}(x) + \beta \cdot z_D(x), -2.5, +2.5 \right)
\end{equation}

\section{System Architecture \& Algorithms}

\subsection{Zero-Cost Probing via Greedy Descent}
During canonical HNSW insertion, vector $x$ traverses the graph from top layer $L_{\text{max}}$ down to the target insertion level $l_{\text{target}}$ using greedy 1-NN search. At level $l_{\text{target}}$, AdaptiveVec performs a single probe beam search with candidate budget $k_{\text{probe}} = \min(25, efC / 4)$. The resulting candidate distances directly supply the input to the LID and Density estimators, incurring less than $0.8\%$ computational overhead.

\subsection{Dynamic Degree \& Construction Allocation}
AdaptiveVec maps the difficulty score $S(x)$ to an adaptive base edge quota $M(x)$ and construction beam $efConstruction(x)$:
\begin{align}
M(x) &= \text{clip}\left( \lfloor M_{\text{base}} \cdot (1 + \gamma \cdot S(x)) \rceil, M_{\text{min}}, M_{\text{max}} \right) \\
efC(x) &= \text{clip}\left( \lfloor efC_{\text{base}} \cdot (1 + \gamma \cdot S(x)) \rceil, efC_{\text{min}}, efC_{\text{max}} \right)
\end{align}

\subsection{Layer-Decoupled Link Scaling}
For higher layers $l \ge 1$, which act as coarse routing expressways, edge capacity is geometrically compressed:
\begin{equation}
M^{(l)}(x) = \max\left( M_{\text{min}}^{(l)}, \lfloor M(x) \cdot \lambda^l \rfloor \right)
\end{equation}
where decay factor $\lambda = 0.75$ and $M_{\text{min}}^{(l)} = 4$.

\subsection{Hubness-Aware In-Degree Regulation}
During Relative Neighborhood Graph (RNG) candidate evaluation, metric distance is regularized by in-degree centrality:
\begin{equation}
d_{\text{eff}}(u, v) = d(u, v) \cdot \left[ 1 + \mu \cdot \left( \frac{\text{deg}_{\text{in}}(v)}{\bar{\text{deg}}_{\text{in}}} \right) \right]
\end{equation}
where $\mu = 0.15$. This penalizes over-connected hubs and redirects links to topologically diverse alternatives.

\subsection{Distance Stagnation Early Exit}
During query-time beam search at layer 0, standard HNSW evaluates candidate nodes until the priority queue is exhausted. AdaptiveVec tracks the minimum discovered candidate distance $d_{\text{best}}$. If absolute improvement $\Delta = d_{\text{best}}^{(\text{old})} - d_{\text{best}}^{(\text{new})}$ satisfies $\Delta < \epsilon$ ($\epsilon = 10^{-4}$) for $p = 6$ consecutive hops, search terminates immediately. This introduces a controlled trade-off: an intentional 1.68\% recall delta ($0.9913 \to 0.9745$) eliminates 30.1\% of distance evaluations ($1,121.2 \to 783.5$ evals/query) and boosts throughput from 4,708.1 to 7,075.3 QPS (+50.3\%).

\textit{Remark on Implementation Correctness (Min-Heap Monotonicity Bug):} In early prototyping, $d_{\text{best}}$ was evaluated against individual popped candidate distances from the traversal priority queue. Because popped elements from a min-heap are non-decreasing, evaluating consecutive differences across popped elements caused false-positive stagnation triggers prematurely during early exploratory jumps, devastating recall. The resolution enforced tracking the strictly non-increasing global best distance across the entire visited candidate set ($d_{\text{best}} \leftarrow \min(d_{\text{best}}, d_{\text{new}})$).

\section{Theoretical Complexity Analysis}

\textbf{Proposition 1 (Expected Memory Footprint Complexity).} Let $N$ be the total vector count. In canonical HNSW with fixed degree quota $M$, graph edge memory is $E_{\text{HNSW}} = 2 \cdot M \cdot N \cdot \text{sizeof}(\text{uint32})$. In AdaptiveVec, edge memory scales as:
\begin{equation}
E_{\text{AdaptiveVec}} = 2 \cdot N \cdot \mathbb{E}[M(x)] \cdot \text{sizeof}(\text{uint32})
\end{equation}
Under clustered, low-LID subspaces where $\mathbb{E}[S(x)] < 0$, empirical expectation satisfies $\mathbb{E}[M(x)] < M_{\text{base}}$. Combined with layer-decoupled scaling, this delivers a measured 7.4\% edge reduction on SIFT-100K while maintaining graph connectivity. With Asymmetric INT8 quantization (SQ8), index RAM drops by 62.4\% (59.9MB to 22.5MB).

\textbf{Proposition 2 (Insertion Complexity Bounds).} With $k_{\text{probe}} \le 25 \ll efC_{\text{base}}$, the signal probing term is $\mathcal{O}(k_{\text{probe}} \cdot D + k_{\text{probe}} \log k_{\text{probe}})$, strictly dominated by layer construction search. Discounted construction budgets in low-LID clusters accelerate total indexing time by 28.1\% (33.5s vs. 46.6s baseline).

\section{Empirical Evaluation}

\subsection{Experimental Setup}
Evaluations were conducted on a single-node host: Intel Core 5 210H (8 cores / 12 threads: 4 Performance cores up to 4.8 GHz, 4 Efficient cores, 12MB L3 cache), 16.0 GB DDR5 RAM, Windows 11 Home Single Language (64-bit, build 26100), utilizing MSYS2 MinGW-w64 `g++ 16.1.0` with flags `-O3 -mavx2 -mfma -std=c++17`. Python benchmarks used Python 3.11 with NumPy 1.26.

\textbf{Evaluated Corpora:}
(1) \textit{SIFT-100K Subset (Texmex IRISA):} 100,000 vectors ($D=128$, $L_2$ space, verified MD5 \texttt{b23d1b3b2ee8469d819b61ca900ef0ed}), evaluated against 10,000 genuine test queries with exact brute-force ground truth.
(2) \textit{Synthetic-Multi-Cluster:} 50,000 vectors across 8 Gaussian clusters ($D=64$, $L_2$ space, 1,000 test queries).

\textbf{Corpora Explicitly Not Evaluated:}
(1) \textit{DBpedia-100K:} Not evaluated due to the lack of verified raw OpenAI \texttt{text-embedding-3-small} embeddings locally; synthetic proxies were strictly excluded.
(2) \textit{Texmex SIFT-1M (Full) \& Stanford GloVe-100:} Excluded from continuous automated benchmark sweeps due to RAM and execution time budgets on a 16GB RAM host; reserved for future out-of-core scaling sweeps.

\begin{table}[htbp]
\caption{Macro-Benchmark Comparison on SIFT-100K (Intel Core 5 210H)}
\centering
\begin{tabular}{lcccccc}
\toprule
\textbf{Configuration} & \textbf{Edges} & \textbf{Build (s)} & \textbf{RAM} & \textbf{Recall@10} & \textbf{QPS} & \textbf{Evals/q} \\
\midrule
Baseline HNSW & 2,709,125 & 46.6 & 59.9 MB & \textbf{0.9913} & 4,708.1 & 1,121.2 \\
Regime A (Step 5) & 2,509,138 & \textbf{33.5} & 59.2 MB (Parity) & 0.9745 & \textbf{7,075.3} & \textbf{783.5} \\
Regime B (Step 6 SQ8) & 2,509,743 & 64.0 & \textbf{22.5 MB} & 0.9594 & 2,987.6 & 842.6 \\
\bottomrule
\end{tabular}
\end{table}

\begin{table}[htbp]
\caption{Controlled Component Ablation Study}
\centering
\begin{tabular}{lccccc}
\toprule
\textbf{Ablation Step} & \textbf{Edges} & \textbf{Build (s)} & \textbf{RAM} & \textbf{Recall} & \textbf{QPS} \\
\midrule
\multicolumn{6}{l}{\textbf{Part A: SIFT-100K Subset ($N=100\text{K}, D=128, Q=10\text{K}$)}} \\
1. Baseline HNSW & 2,709,125 & 46.6 & 59.9MB & 0.9913 & 4,708.1 \\
2. + Dynamic $M(x)$ & 2,549,825 & 47.9 & 59.3MB & 0.9883 & 5,147.1 \\
3. + Layer Scaling & 2,515,277 & 33.5 & 59.2MB & 0.9887 & 2,242.3 \\
4. + Hubness Reg. & 2,510,334 & 35.8 & 59.2MB & 0.9854 & 5,452.0 \\
5. + Stagnation Exit & 2,509,138 & 33.5 & 59.2MB & 0.9745 & \textbf{7,075.3} \\
6. + Asym. SQ8 & 2,509,743 & 64.0 & \textbf{22.5MB} & 0.9594 & 2,987.6 \\
\midrule
\multicolumn{6}{l}{\textbf{Part B: Synthetic-Multi-Cluster ($N=50\text{K}, D=64, Q=1\text{K}$)}} \\
1. Baseline HNSW & 1,284,614 & 28.9 & 17.5MB & 0.9220 & 5,920.7 \\
2. + Dynamic $M(x)$ & 1,288,175 & 15.2 & 17.5MB & 0.9172 & 5,893.3 \\
3. + Layer Scaling & 1,257,271 & 14.8 & 17.4MB & 0.9145 & 6,801.5 \\
4. + Hubness Reg. & 1,289,398 & 15.1 & 17.5MB & 0.7821 & 6,529.0 \\
5. + Stagnation Exit & 1,263,970 & 14.3 & 17.4MB & 0.8566 & \textbf{7,420.7} \\
6. + Asym. SQ8 & 1,293,780 & 14.9 & \textbf{8.4MB} & 0.8564 & 6,610.5 \\
\bottomrule
\end{tabular}
\end{table}

\subsection{Methodology \& AI-Assisted Engineering Disclosure}
In adherence to modern transparency standards, the author discloses that AI-assisted engineering tools (Anthropic Claude and Google DeepMind Antigravity) were utilized for exploratory scaffolding, generating unit tests, and optimizing C++ SIMD AVX2 intrinsic loops. All mathematical formulations, bug resolutions, and empirical evaluations were executed and verified on local hardware testbeds.

\section{Conclusion}
AdaptiveVec demonstrates that uniform hyper-parameterization in proximity graph indexing is fundamentally sub-optimal. By harvesting online geometric signals directly from greedy descent routing paths at $<0.8\%$ overhead, AdaptiveVec dynamically provisions edge quotas, compresses higher-layer links, regulates hubness, and accelerates search termination. Empirical evaluations on canonical SIFT-100K demonstrate 7.4\% edge reductions, +50.3\% query throughput gains (7,075.3 QPS vs 4,708.1 QPS), and 62.4\% total index memory savings (22.5MB vs 59.9MB) with Asymmetric INT8 scalar quantization, establishing a defensible foundation for vector retrieval on commodity hardware.

\begin{thebibliography}{10}
\bibitem{malkov2020hnsw} Y. A. Malkov and D. A. Yashunin, ``Efficient and robust approximate nearest neighbor search using Hierarchical Navigable Small World graphs,'' \textit{IEEE TPAMI}, vol. 42, no. 4, pp. 824--836, 2020.
\bibitem{levina2005mle} E. Levina and P. J. Bickel, ``Maximum likelihood estimation of intrinsic dimension,'' \textit{NeurIPS}, 2005.
\bibitem{amsaleg2015lid} L. Amsaleg et al., ``Estimating local intrinsic dimensionality,'' \textit{ACM SIGKDD}, 2015.
\bibitem{radovanovic2010hubs} M. Radovanović et al., ``Hubs in space: Popular nearest neighbors in high-dimensional data,'' \textit{JMLR}, vol. 11, pp. 2487--2531, 2010.
\bibitem{subramanya2019diskann} S. Subramanya et al., ``DiskANN: Fast accurate billion-point nearest neighbor search on a single node,'' \textit{NeurIPS}, 2019.
\bibitem{fu2019nsg} C. Fu et al., ``Fast approximate nearest neighbor search with the navigating spreading-out graph,'' \textit{PVLDB}, 2019.
\bibitem{johnson2021faiss} J. Johnson et al., ``Billion-scale similarity search with GPUs,'' \textit{IEEE TBD}, 2021.
\bibitem{welford1962note} B. P. Welford, ``Note on a method for calculating corrected sums of squares and products,'' \textit{Technometrics}, 1962.
\bibitem{wang2021survey} M. Wang, X. Xu, Q. Yue, and Y. Wang, ``A comprehensive survey and experimental comparison of graph-based approximate nearest neighbor search,'' \textit{PVLDB}, vol. 14, no. 11, pp. 1964--1978, 2021.
\bibitem{li2020survey} W. Li et al., ``Approximate nearest neighbor search on high dimensional data---experiments, analyses, and improvement,'' \textit{IEEE TKDE}, vol. 32, no. 8, pp. 1475--1488, 2020.
\end{thebibliography}

\end{document}
